% Размер бумаги А4, шрифт 12 пунктов
\documentclass[a4paper,12pt]{report}
% Перекодировка шрифтов
\usepackage[T2A]{fontenc}
% Кодировка: utf8
\usepackage[utf8]{inputenc}
% Используем русский и английский языки с переносами
\usepackage[english,russian]{babel}
% Подключаем нужные пакеты расширений
\usepackage{amssymb,amsfonts,amsmath,mathtext,cite,enumerate,float}

% Вставка рисунков
\usepackage{graphicx}
% Путь к рисункам
\graphicspath{{../pictures/}}

% Рыба
\usepackage{lipsum}
% Подсчёт объектов
\usepackage{totcount}
\regtotcounter{page}
\regtotcounter{figure}
\regtotcounter{table}

% Изменение размеров шрифтов
\usepackage{scrextend}
% Размер по умолчанию — 14 пунктов, интервал полуторный (14*1.25)
\changefontsizes[17.5pt]{14pt}
% Отступ первой строки параграфа после начала раздела
\usepackage{indentfirst}
% Отступ начала параграфа
\setlength{\parindent}{1cm}

% Заменяем библиографию с квадратных скобок на точку
\makeatletter
\renewcommand{\@biblabel}[1]{#1.}
\makeatother

% Геометрия страницы
\usepackage{geometry}
% Поля
\geometry{left=3cm}
\geometry{right=1cm}
\geometry{top=2cm}
\geometry{bottom=2cm}

% Меняем везде перечисления на цифра.цифра
\renewcommand{\theenumi}{\arabic{enumi}}
\renewcommand{\labelenumi}{\arabic{enumi}}
\renewcommand{\theenumii}{.\arabic{enumii}}
\renewcommand{\labelenumii}{\arabic{enumi}.\arabic{enumii}.}
\renewcommand{\theenumiii}{.\arabic{enumiii}}
\renewcommand{\labelenumiii}{\arabic{enumi}.\arabic{enumii}.\arabic{enumiii}.}

% Изменяем \part и \chapter из report.cls чтобы они не выводили нумерованные заголовки
% Рецепт отсюда: http://tex.stackexchange.com/questions/70621/part-and-chapter-commands-output-both-numbered-and-named-headings/70624
\makeatletter
    \def\@part[#1]#2{%
        \ifnum \c@secnumdepth >-2\relax
        \refstepcounter{part}%
        \addcontentsline{toc}{part}{#1}%
        \else
        \addcontentsline{toc}{part}{#1}%
        \fi
        \markboth{}{}%
        {
            \centering
            \interlinepenalty \@M
            \normalfont
            \Huge \bfseries #2\par}%
        \@endpart}
    \def\@chapter[#1]#2{\ifnum \c@secnumdepth >\m@ne
        \refstepcounter{chapter}%
        \typeout{\@chapapp\space\thechapter.}%
        \addcontentsline{toc}{chapter}%
        {#1}%
        \else
        \addcontentsline{toc}{chapter}{#1}%
        \fi
        \chaptermark{#1}%
        \addtocontents{lof}{\protect\addvspace{10\p@}}%
        \addtocontents{lot}{\protect\addvspace{10\p@}}%
        \if@twocolumn
        \@topnewpage[\@makechapterhead{#2}]%
        \else
        \@makechapterhead{#2}%
        \@afterheading
        \fi}
    \def\@makechapterhead#1{%
        \vspace*{50\p@}%
        {\parindent \z@ \raggedright \normalfont
            \interlinepenalty\@M
            \Huge \bfseries #1\par\nobreak
            \vskip 40\p@
    }}
    \renewcommand{\thesection}{\arabic{section}} % Удаляем \thechapter из \thesection
\makeatother

% Для веб-адресов в библиографии
\usepackage{url}
% Библиография в порядке упоминания в тексте
\bibliographystyle{unsrt}

\begin{document}
% Титульный лист
\begin{titlepage}
    \newgeometry{top=2cm,bottom=2cm,left=2cm,right=2cm}
    \newpage
    
    \begin{center}
        Федеральное агенство по образованию Российской Федерации \\
        Московский Государственный Технический Университет \\*
        имени Н.Э.Баумана \\*
        \vspace{-12mm}
        \begin{figure}[h]
            \center{\includegraphics[width=0.2\linewidth]{symbol}}
        \end{figure}
        \vspace{-16mm}
        \hrulefill
    \end{center}
    \center{Факультет «Робототехника и комплексная автоматизация»\\
    		Кафедра «Системы Автоматизированного Проектирования»}
    \begin{center}
        \Large Отчёт о преддипломной практике по теме:
    \end{center}
    
    \vspace{2.5em}
    
    \begin{center}
        \textsc{\textbf{Инструментарий коллективной оптимизации программ}}
    \end{center}
    
    \vspace{6em}
    
    \begin{flushleft}
        \hspace{8.5cm}Исполнитель \hrulefill Панков М.К. \\
        \vspace{1.5em}
        \hspace{8.5cm}Руководитель \hrulefill Карпенко А.П.\\
    \end{flushleft}
    
    \vspace{\fill}
    
    \begin{center}
        Москва, \\*
        2013
    \end{center}

\end{titlepage}
% Аннотация
\section*{Аннотация}
Отчёт о преддипломной технологической практике на тему «Коллективная оптимизация программ» содержит \total{page} страницы машинописного текста, \total{figure} рисунок и \total{table} таблиц. 
% Описание проделанной работы
{}
% Оглавление
\tableofcontents
\section{Введение}
\subsection{Мотивация}
Ниже приведены основные причины заинтересованности в данной задаче.
\begin{itemize}
    \item Многие исследования в области применения машинного обучения в компиляции \cite{Dubach2009PCO,Thomson2009RTT,Stephenson2003MOI} показывают большой потенциал и востребованность решения данной задачи, но плохо систематизированы.
    \item Инженерная задача по созданию инфраструктуры для коллективной оптимизации не решена до сих пор \cite{Fursin2010COP}.
    \item Конечным пользователям компиляторов мешает именно отсутствие инструментария, который помог бы им пользоваться обучаемым компилятором так же, как обычным (используя интерфейс «чёрного ящика»).
    \item Производители аппаратного и программного обеспечения, такие как Intel, IBM, ARM заинтересованы в решении такой задачи \cite{IntelExascale}.
\end{itemize}


\subsection{Цель}
Конечная цель, которой требуется достичь --- разработать систему сбора, систематизации, формализации данных о производительности компилируемых программ в зависимости от настроек компилятора, а также выполняющую функции поддержки базы знаний и обучения. Цели данной системы представлены ниже.
\begin{itemize}
    \item Предсказание производительности программ.
    \item Оптимизация производительности программ путём изменения настроек компилятора.
\end{itemize}

На данном, начальном, этапе работы рассматривается только один критерий оптимальности --- производительность, а именно --- полное время выполнения программы.

Однако, сравнивая стоящую перед нами задачу с текущими успехами в этой области исследований \cite{Fursin2010COP}, обнаруживаем, что достижение данной цели в рамках курсового проекта не представляется возможным. В связи с этим разработка ограничена учебным заданием и продолжение работы будет осуществляться в рамках дипломного проекта.

Итак, цель данной работы --- разработать основные компоненты расширяемого инструментария для проведения экспериментов по оптимизации программ (в ручном и автоматическом режиме) и прозрачного использования данных из экспериментов в промышленной компиляции.


\subsection{Обзор предшествующих работ}
Как было сказано выше, многие исследования в области применения машинного обучения в компиляции \cite{Dubach2009PCO,Thomson2009RTT,Stephenson2003MOI} показывают большой потенциал и востребованность решения данной задачи, но плохо систематизированы. Вследствие этого основной проблемой является реализация общей базы данных программ и экспериментов и каркаса инструментария для проведения исследований. Реализация всех конечных сценариев, необходимых в исследовании, не является приоритетной задачей, но должно быть осуществимо достаточно простым способом силами самих программных инженеров-конечных пользователей системы и других исследователей.

Ниже приведены некоторые отличия данной работы от существующих.
\begin{itemize}
    \item Исследователям и инженерам нужен в первую очередь минимальный функционал для решения описываемой задачи. Графический интерфейс не является необходимым. В сосредоточенности на традиционном текстовом командном интерфейсе основное отличие этой работы от работы \cite{Fursin2010COP}.
    \item Инструментарий не может реализовывать все необходимые сценарии работы с самого начала. Поэтому цель --- сделать систему компактной, модульной и документированной, что обеспечит хорошую поддерживаемость и расширяемость. В этом также одно из отличий от работы \cite{Fursin2010COP}.
    \item Разработка системы сфокусирована на применении локально в изоляции от других исследователей, поскольку это --- первоначально необходимый сценарий использования. Разработка многопользовательских сценариев является следующим шагом. Однако используемая технологическая база позволяет относительно просто расширить инструментарий для использования в распределённых системах.
\end{itemize}


\subsection{Задачи}
Задачи, решаемые в данной работе, перечислены далее в примерном порядке реализации.
\begin{enumerate}
    \item Установить репозиторий исходных кодов программ.
    \item Создать и настроить базу данных для хранения экспериментов.
    \item Реализовать интерактивное текстовое окружение для проведения экспериментов, обработки и сохранения данных.
    \item Реализовать построение графиков для визуального представления данных экспериментов.
\end{enumerate}
% Название библиографии
\renewcommand\bibname{Список использованных источников}
% Список использованных источников
\bibliography{biblio/practice}
\end{document}